\color{unidarkgreen}{\section[(рекомендуемое) Перечень оборудования и средств измерения]{(рекомендуемое)\\Перечень оборудования и средств измерения}\label{app:unit}}
\color{black}

\setlength{\extrarowheight}{0.1cm}
\small
\begin{longtable}{
  |>{\raggedright\arraybackslash}m{7cm}
  |>{\centering\arraybackslash}m{3cm}
  |>{\centering\arraybackslash}m{3cm}
  |>{\centering\arraybackslash}m{3cm}|
}
\caption{Перечень оборудования и средств измерения\hfill\vspace{-0.5\baselineskip}}\label{appunit:tbl1}\\ 
\hline
\rowcolor{gray!30}
\multicolumn{1}{|>{\centering\arraybackslash}m{7cm}|}{Наименование оборудования} & 
Диапазон измеряемых (контролируемых) величин & 
Класс точности или погрешность измерения & 
Рекомендованное оборудование или нормативный документ \\
\hline
\endfirsthead

\caption*{\hspace{3pt}\emph{Продолжение таблицы \ref{appunit:tbl1}\hfill\vspace{-0.5\baselineskip}}} \\
\hline
\rowcolor{gray!30}
\multicolumn{1}{|>{\centering\arraybackslash}m{7cm}|}{Наименование оборудования} & 
Диапазон измеряемых (контролируемых) величин & 
Класс точности или погрешность измерения & 
Рекомендованное оборудование или нормативный документ \\ 
\hline
\endhead
\hline
\endfoot
\endlastfoot

% Тело таблицы:
Вольтметр переменного тока & до 300 В & 0,5 & ГОСТ 8711-93 \\ \hline
Вольтметр постоянного тока & до 300 В & 0,5 & ГОСТ 8711-93 \\ \hline
Амперметр переменного тока & (2,5-5) А & 0,5 & ГОСТ 8711-93 \\ \hline
Мультиметр цифровой & \makecell[c]{(0-750) В \\ (0-10) А} & ±0,1\% & \makecell[c]{APPA-107N \\ APPA-109N} \\ \hline
Источник питания постоянного тока & \makecell[c]{(0-6) А \\ (0-30) В} & ±0,5\% & \makecell[c]{GPS-2303 \\ GPS-3303 \\ GPS-4303} \\ \hline
Источник питания постоянного тока & \makecell[c]{(8-300) В \\ (1-30) А} & \makecell[c]{±(0,005 Uуст+0,2 В) \\ ±(0,005 Iуст+0,02 A)} & GPR \\ \hline
Мегаомметр  & \makecell[c]{(0-1000) МОм \\ 1000 В} & ± 15\% & ГОСТ 23706-93 \\ \hline
Миллиомметр  & от 10 мкОм до 1 кОм & ±(0,1\%+0,5 мкОм) & МИКО-7 \\ \hline
Электронный осциллограф & \makecell[c]{(0-300) В \\ (5-400) Гц} & \makecell[c]{± 10\% \\ ± 1\%} & ГОСТ 9829-81 \\ \hline
Измеритель временных параметров реле & (0-100) с & 0,005/0,004 & ТУ25-04-08.003-83 \\ \hline
Штангенциркуль & (0-400) мм & ±0,05 мм & \makecell[c]{ШЦ-II-400-0,05 \\ ГОСТ 166} \\ \hline
Измеритель сопротивления, увлажненности и степени старения электроизоляции & \makecell[c]{(50-2500) В, 50 Гц} & \makecell[c]{± 10\%} & \makecell[c]{MIC-2500} \\ \hline
Установка для проверки электрической безопасности & \makecell[c]{(100-5000) В} & \makecell[c]{± (0,03 U+30 В)} & \makecell[c]{GPT-815} \\ \hline
\multicolumn{4}{|p{15cm}|}{\textmd{Примечание – При проведении испытаний и проверок допускается применение другого оборудования, обеспечивающего измерение контролируемых параметров с точностью не ниже требуемой.}} \\ \hline

\end{longtable}
\normalsize
\setlength{\extrarowheight}{0.0cm}



\clearpage